% QUESTAO
A tabela de frequências a seguir apresenta o levantamento salarial dos funcionários de uma empresa, com os valores agrupados em quatro intervalos.

\renewcommand{\arraystretch}{1.3} % Espaço extra para o aluno escrever os cálculos
\arrayrulecolor{CorSerie}    % Linhas na cor do tema
\noindent\begin{tabular}{|C{28mm}|C{9mm}|C{9mm}|}
\hline
\rowcolor{CorSerie!70} 
\textbf{\color{white}Salário (R\$)} & 
\textbf{\color{white}FA} & 
\textbf{\color{white}FR} \\ \hline
% Classe 1
\rowcolor{CorSerie!5}
\rule{8mm}{0pt} \tikz[thick, CorSerie]{\draw (0,-0.15)--(0,0.15) (0,0)--(0.7,0);} \rule{8mm}{0pt} & & 40\% \\ \hline
% Classe 2
\rule{8mm}{0pt} \tikz[thick, CorSerie]{\draw (0,-0.15)--(0,0.15) (0,0)--(0.7,0);} \rule{8mm}{0pt} & 15 & \\ \hline
% Classe 3
\rowcolor{CorSerie!5}
2200 \tikz[thick, CorSerie]{\draw (0,-0.15)--(0,0.15) (0,0)--(0.7,0);} \rule{8mm}{0pt} &60 &30\%\\\hline
% Classe 4
\rule{8mm}{0pt} \tikz[thick, CorSerie]{\draw (0,-0.15)--(0,0.15) (0,0)--(0.7,0);} 3100 &&\\\hline
% Total
\rowcolor{CorSerie!15}
\textbf{Total} & &  \\ \hline
\end{tabular}

Preencha os campos em branco na tabela e determine: qual é o total de funcionários da empresa e qual é a amplitude dos intervalos de classe?

% ALTERNATIVAS
total $=200$ e amplitude $=450$
total $=90$ e amplitude $=100$
total $=140$ e amplitude $=200$
total $=70$ e amplitude $=250$
total $=110$ e amplitude $=600$
